\subsection{全体概要：この章で学ぶこと}
第17章は、ソフトウェア技術者が「曖昧でない考え方」をコード、仕様、モデル、テストへ落とし込むための数学を扱う。ここでいう数学は、計算練習としての算術だけではない。中心にあるのは、論理、集合、関係、関数、グラフ、状態、文法、確率、誤差、代数構造といった、ソフトウェアを正確に考えるための道具である。

この章の狙いは、複雑な対象を形式化し、推論し、妥当な結論を得るための基礎を理解することにある。たとえば、if文の条件、データ構造、状態遷移、パーサ、暗号、乱択アルゴリズム、数値シミュレーションは、いずれも数学的な概念に支えられている。

\begin{notebox}{この資料の主張}
数学的基礎は、ソフトウェア開発における「正しく考えるための共通語」である。論理は条件分岐と仕様の意味を支え、集合と関数はデータと処理の関係を支え、グラフと木は構造を支え、有限状態機械と文法は振る舞いと言語処理を支える。確率と誤差は、不確実性と近似を扱うために必要である。
\end{notebox}
到達目標 この章を読み終えると、次のことを説明できるようになる。

\begin{itemize}[leftmargin=2em]
\item 命題論理と述語論理の違いを説明できる。
\item 直接証明、背理法、数学的帰納法、例による存在証明を区別できる。
\item 集合、関係、関数をデータの整理方法として説明できる。
\item グラフ、木、二分探索木、有限状態機械がソフトウェアで使われる理由を説明できる。
\item 文法、正規表現、文脈自由文法がプログラミング言語や入力処理と関係する理由を説明できる。
\item 数え上げ、離散確率、数値誤差を、テスト、性能評価、シミュレーションに結び付けて説明できる。
\end{itemize}
学習順序 初学者は、まず論理と集合を読むとよい。これらは後のグラフ、木、状態機械、文法、確率の土台になる。次に、データ構造に近いグラフと木、振る舞いに近い有限状態機械と文法を読む。最後に、数論、数え上げ、確率、数値誤差、代数構造、微積分を、用途と結び付けて読むとよい。

\par\smallskip
\begin{center}
\centering
\IfFileExists{assets/generated_figures/ch17/fig-ch17-01.png}{\includegraphics[width=0.96\linewidth]{assets/generated_figures/ch17/fig-ch17-01.png}}{\fbox{\parbox[c][48mm][c]{0.9\linewidth}{\centering 画像生成待ち\\17章の全体地図}}}
\par\smallskip
\refstepcounter{figure}\label{fig:ch17-01}図 \thefigure: 17章の全体地図
\end{center}
図 \thefigure は、17章の全体地図を視覚的に整理したものである。
\par\smallskip
\subsubsection{最初に必要な用語}
\par\smallskip\noindent\refstepcounter{table}\label{tab:ch17-01}表 \thetable: 最初に必要な用語における用語・初学者向けの説明\par\smallskip
表 \thetable は、最初に必要な用語における用語・初学者向けの説明を比較・確認しやすい形で整理したものである。\par\smallskip
\begin{longtable}{>{\raggedright\arraybackslash}p{0.26\linewidth} >{\raggedright\arraybackslash}p{0.68\linewidth}}
\toprule
用語 & 初学者向けの説明 \\
\midrule
\endfirsthead
\toprule
用語 & 初学者向けの説明 \\
\midrule
\endhead
形式的 & 意味が曖昧にならないよう、記号、規則、推論手順を明確に定めること。 \\
\term{命題}{Proposition} & 真または偽のどちらかに判定できる文である。 \\
\term{述語}{Predicate} & 対象について性質や関係を述べる文の型である。「xは赤い」「xはyより大きい」などである。 \\
\term{証明}{Proof} & ある主張が必ず真であることを、認めた規則から厳密に示す説明である。 \\
\term{集合}{Set} & 対象をひとまとまりにしたものである。要素の順序や重複ではなく、含まれるかどうかに注目する。 \\
関係 & 複数の集合の要素どうしの対応である。例として、人と電話番号の対応がある。 \\
\term{関数}{Function} & 入力ごとに出力がただ一つ決まる関係である。 \\
\term{グラフ}{Graph} & 点と線で対象同士のつながりを表す構造である。 \\
\term{木}{Tree} & 根から枝分かれする階層構造であり、単純な閉路を持たない。 \\
\term{有限状態機械}{Finite State Machine} & 有限個の状態を持ち、入力によって状態を遷移する抽象モデルである。 \\
文法 & 文字列や文を作る規則である。プログラミング言語や設定ファイルの構文を扱う基礎である。 \\
離散確率 & 数え上げられる結果に確率を割り当てる考え方である。 \\
数値誤差 & 計算機で表した近似値と本当の値との差である。 \\
\bottomrule
\end{longtable}
